% 第 4 章 基于指针的图书馆管理系统
\renewcommand{\codeInset}{18bp}
本章我们将继续研究图书馆管理系统。与第 3 章「构建图书馆管理系统」一样，我们同样有三个类：\texttt{Book}、\texttt{Customer} 和 \texttt{Library}。不过，这一次有一个重大区别：我们不再使用编号，而是使用指针。每个 \texttt{Book} 对象持有一个指针，指向借走这本书的顾客（\texttt{Customer} 类的对象），同时还持有一份指针列表，指向预约这本书的各位顾客。同样，每位顾客也持有多组指针，指向自己借阅和预约的图书（\texttt{Book} 类的对象）。

不过，这种做法也带来了一个问题：我们无法把指针的值直接存入文件。因此，保存文件时需要把指针转换成图书列表和顾客列表中的索引，而加载文件时则需要把索引再转换回指针。这一过程称为 marshmallowing（直译为「棉花糖化」）。

本章将深入讨论以下主题：

\begin{itemize}
\item 与第 3 章「构建图书馆管理系统」一样，我们仍将使用构成一个小型数据库的图书类和顾客类。不过本章将直接使用指针，而不再使用整数编号。
\item 由于使用指针而不是整数编号，文件处理变得更加复杂。我们需要执行一个称为 marshmallowing 的过程。
\item 最后，我们将使用 C++ 标准库中的泛型类 \texttt{set} 和 \texttt{list}。不过本章中，它们存放的是指向图书对象和顾客对象的指针，而不是对象本身。
\end{itemize}

\section{Book 类}

与上一章的系统类似，我们同样有三个类：\texttt{Book}、\texttt{Customer} 和 \texttt{Library}。\texttt{Book} 类负责记录一本书的信息，每本书都有作者和书名。\texttt{Customer} 类负责记录一位顾客的信息，每位顾客都有姓名和地址。\texttt{Library} 类负责图书馆的各项操作，如借阅、归还和预约。最后，main 函数只是创建一个 \texttt{Library} 类的对象。

\texttt{Book} 类与第 3 章「构建图书馆管理系统」中的 \texttt{Book} 类相似。唯一实质的区别是：这里没有编号，只有指针。

\medskip\noindent{\ttfamily\bfseries Book.h}

\begin{codebox}
class Customer;

class Book {
  public:
    Book();
    Book(const string& author, const string& title);
    const string& author() const { return m_author; }
    const string& title() const { return m_title; }
    void read(ifstream& inStream);
    void write(ofstream& outStream) const;
    int reserveBook(Customer* customerPtr);
    void removeReservation(Customer* customerPtr);
    void returnBook();
\end{codebox}

我们没有返回图书编号的方法，因为本章的图书不使用编号。

\texttt{borrowerPtr} 方法返回借走这本书的顾客的地址；如果这本书当前无人借阅，则返回 \texttt{nullptr}。它有两个版本：第一个版本返回对顾客对象指针的引用，这样我们就可以为该顾客指针赋予新值；第二个版本是常量版本，意味着它可以在常量对象上调用：

\begin{codebox}
    Customer*& borrowerPtr() { return m_borrowerPtr; }
    const Customer* borrowerPtr() const { return m_borrowerPtr; }
\end{codebox}

注意，本章没有 \texttt{borrowed} 方法。我们不需要它，因为如果这本书当前无人借阅，\texttt{borrowerPtr} 就会返回 \texttt{nullptr}。

本章中，\texttt{reservationPtrList} 返回的是顾客指针列表，而不是整数值。它同样有两个版本：第一个版本返回对列表的引用，这样我们就可以在列表中添加和移除指针；第二个版本是常量版本，返回常量列表，这意味着它可以在常量 \texttt{Book} 对象上调用，并且返回的列表不可修改：

\begin{codebox}
    list<Customer*>& reservationPtrList()
        { return m_reservationPtrList; }
    const list<Customer*> reservationPtrList() const
        { return m_reservationPtrList; }
\end{codebox}

输出流运算符的用法与第 3 章「构建图书馆管理系统」相同：

\begin{codebox}
    friend ostream& operator<<(ostream& outStream,
        const Book& book);
\end{codebox}

\texttt{m\_author} 和 \texttt{m\_title} 字段是字符串，与第 3 章「构建图书馆管理系统」类似：

\begin{codebox}
  private:
    string m_author, m_title;
\end{codebox}

不过，我们省去了 \texttt{m\_bookId} 字段，因为本章不使用编号。我们还将 \texttt{m\_borrowedId} 和 \texttt{m\_customerId} 字段替换为 \texttt{m\_borrowerPtr}，它被初始化为 \texttt{nullptr}，因为这本书从一开始并未被借阅：

\begin{codebox}
    Customer* m_borrowerPtr = nullptr;
\end{codebox}

\texttt{m\_reservationPtrList} 字段存放的是指向预约这本书的各位顾客的指针列表，而不是第 3 章「构建图书馆管理系统」中的整数编号列表：

\begin{codebox}
    list<Customer*> m_reservationPtrList;
};
\end{codebox}

\medskip\noindent{\ttfamily\bfseries Book.cpp}

\begin{codebox}
#include <Set>
#include <Map>
#include <String>
#include <FStream>
#include <Algorithm>
using namespace std;

#include "Book.h"
#include "Customer.h"
#include "Library.h"
\end{codebox}

默认构造函数与第 3 章「构建图书馆管理系统」中的构造函数类似：

\begin{codebox}
Book::Book() {
  // Empty.
}
\end{codebox}

第二个构造函数也与第 3 章「构建图书馆管理系统」中的构造函数类似，不过这里没有需要初始化的 \texttt{m\_bookId} 字段：

\begin{codebox}
Book::Book(const string& author, const string& title)
  :m_author(author),
   m_title(title) {
  // Empty.
}
\end{codebox}

\subsection{读取与写入图书}

\texttt{read} 和 \texttt{write} 方法在本章做了简化：它们只读取和写入书的作者与书名。可能存在的借阅者与预约列表，则由 \texttt{Library} 类的 \texttt{save} 和 \texttt{load} 方法来读取和写入：

\begin{codebox}
void Book::read(ifstream& inStream) {
  getline(inStream, m_author);
  getline(inStream, m_title);
}
\end{codebox}

\begin{codebox}
void Book::write(ofstream& outStream) const {
  outStream << m_author << endl;
  outStream << m_title << endl;
}
\end{codebox}

\subsection{借阅与预约图书}

当顾客预约一本书时，指向该 \texttt{Customer} 对象的指针会被加入这本书的预约指针列表，并返回列表的大小，以便通知顾客其在预约列表中的位置：

\begin{codebox}
int Book::reserveBook(Customer* borrowerPtr) {
  m_reservationPtrList.push_back(borrowerPtr);
  return m_reservationPtrList.size();
}
\end{codebox}

当顾客归还一本书时，我们只需把 \texttt{m\_borrowerPtr} 设为 \texttt{nullptr}，表示这本书已不再被借阅：

\begin{codebox}
void Book::returnBook() {
  m_borrowerPtr = nullptr;
}
\end{codebox}

\texttt{removeReservation} 方法只是把该顾客指针从预约列表中移除：

\begin{codebox}
void Book::removeReservation(Customer* customerPtr) {
  m_reservationPtrList.remove(customerPtr);
}
\end{codebox}

\subsection{显示图书}

输出流运算符写出书名和作者、借走这本书的顾客（如果有），以及预约这本书的各位顾客（如果有）：

\begin{codebox}
ostream& operator<<(ostream& outStream, const Book& book) {
  outStream << "\"" << book.m_title << "\" by "
            << book.m_author;
\end{codebox}

如果这本书已被借阅，我们就把借阅者写入流：

\begin{codebox}
  if (book.m_borrowerPtr != nullptr) {
    outStream << endl << "  Borrowed by: "
              << book.m_borrowerPtr->name() << ".";
  }
\end{codebox}

如果这本书的预约列表非空，我们就遍历该列表，并为每一项预约写出对应的顾客：

\begin{codebox}
  if (!book.m_reservationPtrList.empty()) {
    outStream << endl << "  Reserved by: ";
    bool first = true;
    for (Customer* customerPtr : book.m_reservationPtrList) {
      outStream << (first ? "" : ",") << customerPtr->name();
      first = false;
    }
    outStream << ".";
  }
  return outStream;
\end{codebox}

\section{Customer 类}

本章的 \texttt{Customer} 类与第 3 章「构建图书馆管理系统」中的 \texttt{Customer} 类相似。同样，这里的区别在于我们使用的是指针，而不是整数编号。

\medskip\noindent{\ttfamily\bfseries Customer.h}

\begin{codebox}
class Customer {
  public:
    Customer();
    Customer(const string& name, const string& address);
    const string& name() const { return m_name; }
    const string& address() const { return m_address; }
    void read(ifstream& inStream);
    void write(ofstream& outStream) const;
\end{codebox}

\texttt{borrowBook}、\texttt{returnBook}、\texttt{reserveBook} 和 \texttt{unreserveBook} 方法都以指向 \texttt{Book} 对象的指针作为参数：

\begin{codebox}
    void borrowBook(Book* bookPtr);
    void returnBook(Book* bookPtr);
    void reserveBook(Book* bookPtr);
    void unreserveBook(Book* bookPtr);
\end{codebox}

\texttt{loanPtrSet} 和 \texttt{reservationPtrSet} 方法返回的是 \texttt{Book} 指针的集合，而不是整数编号的集合：

\begin{codebox}
    set<Book*>& loanPtrSet() { return m_loanPtrSet; }
    const set<Book*> loanPtrSet() const { return m_loanPtrSet; }
    set<Book*>& reservationPtrSet()
        { return m_reservationPtrSet; }
    const set<Book*> reservationPtrSet() const
        { return m_reservationPtrSet; }
\end{codebox}

与第 3 章「构建图书馆管理系统」相比，输出流运算符没有变化：

\begin{codebox}
    friend ostream& operator<<(ostream& outStream,
        const Customer& customer);
\end{codebox}

\texttt{m\_name} 和 \texttt{m\_address} 字段存放顾客的姓名和地址，与第 3 章「构建图书馆管理系统」一样：

\begin{codebox}
  private:
    string m_name, m_address;
\end{codebox}

\texttt{m\_loanPtrSet} 和 \texttt{m\_reservationPtrSet} 字段存放的是指向 \texttt{Book} 对象的指针，而不是整数编号：

\begin{codebox}
    set<Book*> m_loanPtrSet, m_reservationPtrSet;
};
\end{codebox}

\medskip\noindent{\ttfamily\bfseries Customer.cpp}

\begin{codebox}
#include <Set>
#include <Map>
#include <String>
#include <FStream>
using namespace std;

#include "Book.h"
#include "Customer.h"
#include "Library.h"
\end{codebox}

构造函数与第 3 章「构建图书馆管理系统」中的构造函数类似。第一个构造函数什么也不做，在从文件加载顾客列表时被调用：

\begin{codebox}
Customer::Customer() {
  // Empty.
}
\end{codebox}

第二个构造函数初始化顾客的姓名和地址。不过，与第 3 章「构建图书馆管理系统」的构造函数相比，这里没有需要初始化的 \texttt{m\_customerId} 字段：

\begin{codebox}
Customer::Customer(const string& name, const string& address)
  :m_name(name),
   m_address(address) {
  // Empty.
}
\end{codebox}

\subsection{读取与写入顾客}

与前面 \texttt{Book} 的情况类似，\texttt{read} 和 \texttt{write} 方法也做了简化：它们只读取和写入姓名与地址。借阅集合与预约集合的读取和写入则放在 \texttt{Library} 类中完成，如下所示：

\begin{codebox}
void Customer::read(ifstream& inStream) {
  getline(inStream, m_name);
  getline(inStream, m_address);
}
\end{codebox}

\begin{codebox}
void Customer::write(ofstream& outStream) const {
  outStream << m_name << endl;
  outStream << m_address << endl;
}
\end{codebox}

\subsection{借阅与预约图书}

\texttt{borrowBook} 方法把图书指针加入借阅集合，如果这本书此前被预约过，则同时把它从预约集合中移除：

\begin{codebox}
void Customer::borrowBook(Book* bookPtr) {
  m_loanPtrSet.insert(bookPtr);
  m_reservationPtrSet.erase(bookPtr);
}
\end{codebox}

\texttt{reserveBook} 方法只是把图书指针加入预约集合；而 \texttt{returnBook} 和 \texttt{unreserveBook} 则分别把图书指针从借阅集合和预约集合中移除：

\begin{codebox}
void Customer::reserveBook(Book* bookPtr) {
  m_reservationPtrSet.insert(bookPtr);
}
\end{codebox}

\begin{codebox}
void Customer::returnBook(Book* bookPtr) {
  m_loanPtrSet.erase(bookPtr);
}

void Customer::unreserveBook(Book* bookPtr) {
  m_reservationPtrSet.erase(bookPtr);
}
\end{codebox}

\subsection{显示顾客}

输出流运算符的用法与第 3 章「构建图书馆管理系统」相同：它写出顾客的姓名和地址，以及借阅和预约的图书集合（如果有）：

\begin{codebox}
ostream& operator<<(ostream& outStream,
                    const Customer& customer) {
  outStream << customer.m_name << ", "
            << customer.m_address << ".";
\end{codebox}

如果该顾客的借阅集合非空，我们就遍历它，并为每一次借阅写出对应的图书：

\begin{codebox}
  if (!customer.m_loanPtrSet.empty()) {
    outStream << endl << "  Borrowed books: ";
    bool first = true;
    for (const Book* bookPtr : customer.m_loanPtrSet) {
      outStream << (first ? "" : ", ") << bookPtr->author();
      first = false;
    }
  }
\end{codebox}

同样，如果该顾客的预约集合非空，我们就遍历它，并为每一项预约写出对应的图书：

\begin{codebox}
  if (!customer.m_reservationPtrSet.empty()) {
    outStream << endl << "  Reserved books: ";
    bool first = true;
    for (Book* bookPtr : customer.m_reservationPtrSet) {
      outStream << (first ? "" : ", ") << bookPtr->author();
      first = false;
    }
  }
  return outStream;
\end{codebox}

\section{Library 类}

\texttt{Library} 类与第 3 章「构建图书馆管理系统」中的 \texttt{Library} 类十分相似。不过，我们新增了一些查找方法，用于在把图书馆信息保存到文件以及从文件加载时，在指针与列表索引之间进行转换：

\medskip\noindent{\ttfamily\bfseries Library.h}

\begin{codebox}
class Library {
  public:
    Library();
\end{codebox}

析构函数释放本应用程序中所有动态分配的内存：

\begin{codebox}
    ~Library();
  private:
    static string s_binaryPath;
\end{codebox}

\texttt{lookupBook} 和 \texttt{lookupCustomer} 方法返回指向 \texttt{Book} 和 \texttt{Customer} 对象的指针。如果图书或顾客不存在，则返回 \texttt{nullptr}：

\begin{codebox}
    Book* lookupBook(const string& author,
                     const string& title);
    Customer* lookupCustomer(const string& name,
                             const string& address);
    void addBook();
    void deleteBook();
    void listBooks();
    void addCustomer();
    void deleteCustomer();
    void listCustomers();
    void borrowBook();
    void reserveBook();
    void returnBook();
\end{codebox}

\texttt{lookupBookIndex} 和 \texttt{lookupCustomerIndex} 方法接受一个指针，在图书列表和顾客列表中查找该指针所指的对象，并返回它在列表中的索引：

\begin{codebox}
    int lookupBookIndex(const Book* bookPtr);
    int lookupCustomerIndex(const Customer* customerPtr);
\end{codebox}

\texttt{lookupBookPtr} 和 \texttt{lookupCustomerPtr} 方法接受一个索引，返回图书列表和顾客列表中相应位置上的对象的指针：

\begin{codebox}
    Book* lookupBookPtr(int bookIndex);
    Customer* lookupCustomerPtr(int customerIndex);
\end{codebox}

\texttt{save} 和 \texttt{load} 方法负责把图书馆信息保存到文件以及从文件加载。不过，它们比第 3 章「构建图书馆管理系统」中的对应方法要复杂得多：

\begin{codebox}
    void save();
    void load();
\end{codebox}

\texttt{m\_bookPtrList} 和 \texttt{m\_customerPtrList} 字段存放的是指向 \texttt{Book} 和 \texttt{Customer} 对象的指针，而不是像第 3 章「构建图书馆管理系统」那样直接存放对象本身：

\begin{codebox}
    list<Book*> m_bookPtrList;
    list<Customer*> m_customerPtrList;
};
\end{codebox}

\medskip\noindent{\ttfamily\bfseries Library.cpp}

\begin{codebox}
#include <Set>
#include <Map>
#include <List>
#include <String>
#include <FStream>
#include <IOStream>
#include <CAssert>
using namespace std;

#include "Book.h"
#include "Customer.h"
#include "Library.h"

string Library::s_binaryPath("C:\\Users\\Stefan\\Library.binary");
\end{codebox}

构造函数与第 3 章「构建图书馆管理系统」中的构造函数完全相同：

\begin{codebox}
Library::Library() {
  load();
  bool quit = false;
  while (!quit) {
    cout << "1. Add Book" << endl
         << "2. Delete Book" << endl
         << "3. List Books" << endl
         << "4. Add Customer" << endl
         << "5. Delete Customer" << endl
         << "6. List Customers" << endl
         << "7. Borrow Book" << endl
         << "8. Reserve Book" << endl
         << "9. Return Book" << endl
         << "0. Quit" << endl
         << ": ";
    int choice;
    cin >> choice;
    cout << endl;
    switch (choice) {
      case 1:
        addBook();
        break;
      case 2:
        deleteBook();
        break;
      case 3:
        listBooks();
        break;
      case 4:
        addCustomer();
        break;
      case 5:
        deleteCustomer();
        break;
      case 6:
        listCustomers();
        break;
      case 7:
        borrowBook();
        break;
      case 8:
        reserveBook();
        break;
      case 9:
        returnBook();
        break;
      case 0:
        quit = true;
        break;
    }
    cout << endl;
  }
  save();
}
\end{codebox}

\subsection{查找图书与顾客}

本章的 \texttt{lookupBook} 方法按作者和书名查找 \texttt{Book} 对象，方式与第 3 章「构建图书馆管理系统」类似。不过，如果找到作者与书名都匹配的 \texttt{Book} 对象，它不会把信息复制到给定对象中，而是直接返回指向该对象的指针。如果没有找到 \texttt{Book} 对象，则返回 \texttt{nullptr}：

\begin{codebox}
Book* Library::lookupBook(const string& author,
                          const string& title) {
  for (Book* bookPtr : m_bookPtrList) {
    if ((bookPtr->author() == author) &&
        (bookPtr->title() == title)) {
      return bookPtr;
    }
  }
  return nullptr;
}
\end{codebox}

同样，\texttt{lookupCustomer} 尝试查找姓名和地址都匹配的 \texttt{Customer} 对象。如果找到该对象，就返回它的指针；如果没有找到，则返回 \texttt{nullptr}：

\begin{codebox}
Customer* Library::lookupCustomer(const string& name,
                                  const string& address) {
  for (Customer* customerPtr : m_customerPtrList) {
    if ((customerPtr->name() == name) &&
        (customerPtr->address() == address)) {
      return customerPtr;
    }
  }
  return nullptr;
}
\end{codebox}

\subsection{新增图书}

\texttt{addBook} 方法提示用户输入作者和书名：

\begin{codebox}
void Library::addBook() {
  string author;
  cout << "Author: ";
  cin >> author;
  string title;
  cout << "Title: ";
  cin >> title;
\end{codebox}

在检查这本书是否已经存在时，我们调用 \texttt{lookupBook}。如果书已存在，就返回指向该 \texttt{Book} 对象的指针；如果不存在，则返回 \texttt{nullptr}。因此，我们测试返回值是否不等于 \texttt{nullptr}。如果不等于 \texttt{nullptr}，说明这本书已经存在，于是显示一条错误消息：

\begin{codebox}
  if (lookupBook(author, title) != nullptr) {
    cout << endl << "The book \"" << title << "\" by "
         << author << " already exists." << endl;
    return;
  }
\end{codebox}

在添加这本书时，我们用 \texttt{new} 运算符动态创建一个新的 \texttt{Book} 对象。我们用 C++ 标准的 \texttt{assert} 宏检查图书指针是否为空。如果指针为空，程序将带着错误信息中止执行：

\begin{codebox}
  Book* bookPtr = new Book(author, title);
  assert(bookPtr != nullptr);
  m_bookPtrList.push_back(bookPtr);
  cout << endl << "Added." << endl;
}
\end{codebox}

\subsection{删除图书}

\texttt{deleteBook} 方法通过提示用户输入书的作者和书名来从图书馆中删除一本书。如果这本书存在，我们就先办理归还、再取消预约，最后删除它：

\begin{codebox}
void Library::deleteBook() {
  string author;
  cout << "Author: ";
  cin >> author;
  string title;
  cout << "Title: ";
  cin >> title;
\end{codebox}

我们通过调用 \texttt{lookupBook} 获得指向该 \texttt{Book} 对象的指针：

\begin{codebox}
  Book* bookPtr = lookupBook(author, title);
\end{codebox}

如果指针为 \texttt{nullptr}，说明这本书不存在，于是显示一条错误消息：

\begin{codebox}
  if (bookPtr == nullptr) {
    cout << endl << "The book \"" << title << "\" by "
         << author << " does not exist." << endl;
    return;
  }
\end{codebox}

我们通过查找借阅者来检查这本书是否已被借出：

\begin{codebox}
  Customer* borrowerPtr = bookPtr->borrowerPtr();
\end{codebox}

如果 \texttt{borrowerPtr} 返回的指针不是 \texttt{nullptr}，我们就调用借阅者的 \texttt{returnBook} 方法归还这本书。这样，这本书就不再登记为该顾客借阅的了：

\begin{codebox}
  if (borrowerPtr != nullptr) {
    borrowerPtr->returnBook(bookPtr);
  }
\end{codebox}

此外，我们还需要检查这本书是否被其他顾客预约过。为此，我们取得这本书的预约列表，并为列表中的每位顾客取消对这本书的预约：

\begin{codebox}
  list<Customer*> reservationPtrList =
      bookPtr->reservationPtrList();
\end{codebox}

注意，我们并不检查该顾客是否真的预约了这本书，而只是简单地取消预约。还要注意，我们不需要把任何对象放回列表，因为我们操作的是指向对象的指针，并不复制对象：

\begin{codebox}
  for (Customer* reserverPtr : reservationPtrList) {
    reserverPtr->unreserveBook(bookPtr);
  }
\end{codebox}

在删除这本书时，我们先从图书指针列表中移除该图书指针，然后释放 \texttt{Book} 对象。先显示消息、再删除图书指针，看起来也许有些奇怪，但顺序必须如此：对象一旦被删除，我们就再也不能对它做任何事情了。如果先删除对象再输出它，会导致内存错误：

\begin{codebox}
  m_bookPtrList.remove(bookPtr);
  cout << endl << "Deleted: " << bookPtr << endl;
  delete bookPtr;
}
\end{codebox}

\subsection{列出图书}

在列出图书时，我们首先检查列表是否为空。如果为空，就只输出 \texttt{No books.}：

\begin{codebox}
void Library::listBooks() {
  if (m_bookPtrList.empty()) {
    cout << "No books." << endl;
    return;
  }
\end{codebox}

如果列表不为空，我们就遍历图书指针列表，对每个图书指针解引用并输出相应的信息：

\begin{codebox}
  for (const Book* bookPtr : m_bookPtrList) {
    cout << (*bookPtr) << endl;
  }
}
\end{codebox}

\subsection{新增顾客}

\texttt{addCustomer} 方法提示用户输入顾客的姓名和地址：

\begin{codebox}
void Library::addCustomer() {
  string name;
  cout << "Name: ";
  cin >> name;
  string address;
  cout << "Address: ";
  cin >> address;
\end{codebox}

如果已有同名同地址的顾客，则显示一条错误消息：

\begin{codebox}
  if (lookupCustomer(name, address) != nullptr) {
    cout << endl << "A customer with name " << name
         << " and address " << address << " already exists."
         << endl;
    return;
  }
\end{codebox}

在添加顾客时，我们动态创建一个新的 \texttt{Customer} 对象，并把它加入顾客对象指针列表：

\begin{codebox}
  Customer* customerPtr = new Customer(name, address);
  assert(customerPtr != nullptr);
  m_customerPtrList.push_back(customerPtr);
  cout << endl << "Added." << endl;
}
\end{codebox}

\subsection{删除顾客}

在删除顾客时，我们先查找该顾客，如果不存在则显示一条错误消息：

\begin{codebox}
void Library::deleteCustomer() {
  string name;
  cout << "Customer name: ";
  cin >> name;
  string address;
  cout << "Address: ";
  cin >> address;
  Customer* customerPtr = lookupCustomer(name, address);
\end{codebox}

如果给定姓名和地址的顾客不存在，则显示一条错误消息。请看下面的代码：

\begin{codebox}
  if (customerPtr == nullptr) {
    cout << endl << "Customer " << name
         << " does not exists." << endl;
    return;
  }
\end{codebox}

如果该顾客至少借阅了一本书，就不能删除该顾客，并显示如下所示的错误消息：

\begin{codebox}
  if (!customerPtr->loanPtrSet().empty()) {
    cout << "The customer " << customerPtr->name()
         << " has borrowed books and cannot be deleted."
         << endl;
    return;
  }
\end{codebox}

不过，如果该顾客没有借阅任何图书，就先把他从图书馆中每本书的预约列表中移除，如下面的代码所示：

\begin{codebox}
  for (Book* bookPtr : m_bookPtrList) {
    bookPtr->removeReservation(customerPtr);
  }
\end{codebox}

然后把该顾客从顾客列表中移除，并用 \texttt{delete} 运算符释放 \texttt{Customer} 对象。再次提醒注意：必须先输出顾客信息，再删除其对象；反过来是行不通的，因为我们无法查看一个已被删除的对象，那样会导致内存错误：

\begin{codebox}
  m_customerPtrList.remove(customerPtr);
  cout << endl << "Deleted." << (*customerPtr) << endl;
  delete customerPtr;
}
\end{codebox}

\subsection{列出顾客}

在列出顾客时，我们遍历顾客列表，并对每个 \texttt{Customer} 对象指针解引用，输出对象的信息：

\begin{codebox}
void Library::listCustomers() {
  if (m_customerPtrList.empty()) {
    cout << "No customers." << endl;
    return;
  }
  for (const Customer* customerPtr: m_customerPtrList) {
    cout << (*customerPtr) << endl;
  }
}
\end{codebox}

\subsection{借阅图书}

在借阅一本书时，我们先提示用户输入作者和书名，如下面的代码片段所示：

\begin{codebox}
void Library::borrowBook() {
  string author;
  cout << "Author: ";
  cin >> author;
  string title;
  cout << "Title: ";
  cin >> title;
\end{codebox}

我们查找这本书，如果不存在，则显示一条错误消息，如下面的代码所示：

\begin{codebox}
  Book* bookPtr = lookupBook(author, title);
  if (bookPtr == nullptr) {
    cout << endl << "There is no book \"" << title
         << "\" by " << author << "." << endl;
    return;
  }
\end{codebox}

如果这本书已被其他顾客借走，就不能再次借阅：

\begin{codebox}
  if (bookPtr->borrowerPtr() != nullptr) {
    cout << endl << "The book \"" << title << "\" by "
         << author << " has already been borrowed." << endl;
    return;
  }
\end{codebox}

我们提示用户输入顾客的姓名和地址：

\begin{codebox}
  string name;
  cout << "Customer name: ";
  cin >> name;
  string address;
  cout << "Address: ";
  cin >> address;
  Customer* customerPtr = lookupCustomer(name, address);
\end{codebox}

如果不存在给定姓名和地址的顾客，则显示一条错误消息：

\begin{codebox}
  if (customerPtr == nullptr) {
    cout << endl << "No customer with name " << name
         << " and address " << address << " exists."
         << endl;
    return;
  }
\end{codebox}

最后，我们把这本书加入该顾客的借阅集合，并把该顾客登记为这本书的借阅者：

\begin{codebox}
  bookPtr->borrowerPtr() = customerPtr;
  customerPtr->borrowBook(bookPtr);
  cout << endl << "Borrowed." << endl;
}
\end{codebox}

\subsection{预约图书}

预约流程与前面的借阅流程类似。我们提示用户输入书的作者和书名，以及顾客的姓名和地址，如下所示：

\begin{codebox}
void Library::reserveBook() {
  string author;
  cout << "Author: ";
  cin >> author;
  string title;
  cout << "Title: ";
  cin >> title;
  Book* bookPtr = lookupBook(author, title);
\end{codebox}

如果这本书不存在，则显示一条错误消息：

\begin{codebox}
  if (bookPtr == nullptr) {
    cout << endl << "There is no book \"" << title
         << "\" by " << author << "." << endl;
    return;
  }
\end{codebox}

如果这本书尚未被借出，就不能预约它。这时我们建议用户直接借阅这本书：

\begin{codebox}
  if (bookPtr->borrowerPtr() == nullptr) {
    cout << endl << "The book \"" << title << "\" by "
         << author << " has not been not borrowed. "
         << "Please borrow the book instead of reserving it."
         << endl;
    return;
  }
\end{codebox}

我们提示用户输入顾客的姓名和地址：

\begin{codebox}
  string name;
  cout << "Customer name: ";
  cin >> name;
  string address;
  cout << "Address: ";
  cin >> address;
  Customer* customerPtr = lookupCustomer(name, address);
\end{codebox}

如果该顾客不存在，则显示一条错误消息：

\begin{codebox}
  if (customerPtr == nullptr) {
    cout << endl << "There is no customer with name "
         << name << " and address " << address << "." << endl;
    return;
  }
\end{codebox}

如果该顾客已经借阅了这本书，就不能再预约它：

\begin{codebox}
  if (bookPtr->borrowerPtr() == customerPtr) {
    cout << endl << "The book has already been borrowed by "
         << name << "." << endl;
    return;
  }
\end{codebox}

最后，我们把该顾客加入这本书的预约列表，并把这本书加入该顾客的预约集合。注意，一本书对应的是一个预约顾客列表，而一位顾客对应的是一个预约图书集合。原因在于：当一本书被归还时，预约列表中的第一位顾客将借走这本书；而顾客的预约集合则没有这样的限制：

\begin{codebox}
  int position = bookPtr->reserveBook(customerPtr);
  customerPtr->reserveBook(bookPtr);
\end{codebox}

我们把该顾客在预约列表中的位置通知给这位顾客：

\begin{codebox}
  cout << endl << position << "nd reserve." << endl;
}
\end{codebox}

\subsection{归还图书}

在归还一本书时，我们提示用户输入作者和书名，但不需要询问是哪位顾客借走了这本书——这一信息已经存放在 \texttt{Book} 对象中：

\begin{codebox}
void Library::returnBook() {
  string author;
  cout << "Author: ";
  cin >> author;
  string title;
  cout << "Title: ";
  cin >> title;
  Book* bookPtr = lookupBook(author, title);
\end{codebox}

如果给定作者和书名的这本书不存在，则显示一条错误消息：

\begin{codebox}
  if (bookPtr == nullptr) {
    cout << endl << "There is no book \"" << title << "\" by "
         << author << "." << endl;
    return;
  }
  Customer* customerPtr = bookPtr->borrowerPtr();
\end{codebox}

如果这本书当前无人借阅，则显示一条错误消息：

\begin{codebox}
  if (customerPtr == nullptr) {
    cout << endl << "The book \"" << title << "\" by "
         << author << " has not been borrowed." << endl;
    return;
  }
  bookPtr->returnBook();
  customerPtr->returnBook(bookPtr);
  cout << endl << "Returned." << endl;
\end{codebox}

归还这本书之后，我们需要查明是否有顾客预约了它：

\begin{codebox}
  list<Customer*>& reservationPtrList =
      bookPtr->reservationPtrList();
\end{codebox}

如果这本书的预约列表中至少有一位顾客，我们就取出排在最前面的顾客，把他从这本书的预约列表中移除，把他登记为这本书的借阅者，并把这本书加入该顾客的借阅集合：

\begin{codebox}
  if (!reservationPtrList.empty()) {
    Customer* newCustomerPtr = reservationPtrList.front();
    reservationPtrList.erase(reservationPtrList.begin());
    bookPtr->borrowBook(newCustomerPtr);
    newCustomerPtr->borrowBook(bookPtr);
    cout << endl << "Borrowed by "
         << newCustomerPtr->name() << endl;
  }
}
\end{codebox}

\subsection{查找图书与顾客（索引与指针）}

在把图书馆信息保存到文件以及从文件加载时，我们需要在指向 \texttt{Book} 和 \texttt{Customer} 对象的指针与图书列表和顾客列表中的索引之间进行转换。\texttt{lookupBookIndex} 方法接受一个指向 \texttt{Book} 对象的指针，并返回它在图书列表中的索引：

\begin{codebox}
int Library::lookupBookIndex(const Book* bookPtr) {
  int index = 0;
  for (Book* testPtr : m_bookPtrList) {
    if (bookPtr == testPtr) {
      return index;
    }
    ++index;
  }
\end{codebox}

如果执行到了这里，\texttt{assert} 宏将带着错误信息中止程序的执行。不过我们本不该到达这一点，因为该 \texttt{Book} 指针应当就在图书指针列表中：

\begin{codebox}
  assert(false);
  return -1;
}
\end{codebox}

\texttt{lookupBookPtr} 方法执行相反的任务：它在图书指针列表中找出 bookIndex 所指位置上的 \texttt{Book} 对象指针。如果索引超出了列表的范围，\texttt{assert} 宏将带着错误信息中止程序的执行。不过这不应该发生，因为所有索引都应当落在列表范围之内：

\begin{codebox}
Book* Library::lookupBookPtr(int bookIndex) {
  assert((bookIndex >= 0) &&
         (bookIndex < ((int) m_bookPtrList.size())));
  auto iterator = m_bookPtrList.begin();
  for (int count = 0; count < bookIndex; ++count) {
    ++iterator;
  }
  return *iterator;
}
\end{codebox}

\texttt{lookupCustomerIndex} 方法以与前面 \texttt{lookupBookIndex} 方法相同的方式，给出 \texttt{Customer} 指针在顾客指针列表中的索引：

\begin{codebox}
int Library::lookupCustomerIndex(const Customer* customerPtr) {
  int index = 0;
  for (Customer* testPtr : m_customerPtrList) {
    if (customerPtr == testPtr) {
      return index;
    }
    ++index;
  }
  assert(false);
  return -1;
}
\end{codebox}

\texttt{lookupCustomerPtr} 方法按照与前面 \texttt{lookupBookPtr} 方法相同的方式，在顾客指针列表中查找给定索引处的 \texttt{Customer} 指针：

\begin{codebox}
Customer* Library::lookupCustomerPtr(int customerIndex) {
  assert((customerIndex >= 0) &&
         (customerIndex < ((int) m_customerPtrList.size())));
  auto iterator = m_customerPtrList.begin();
  for (int count = 0; count < customerIndex; ++count) {
    ++iterator;
  }
  return *iterator;
}
\end{codebox}

\subsection{Marshmallowing（棉花糖化）}

本章 \texttt{Library} 类的 \texttt{save} 和 \texttt{load} 方法比第 3 章「构建图书馆管理系统」中的对应方法要复杂一些。原因在于，我们无法直接保存指针，因为指针存放的内存地址在每次运行之间都会变化。我们必须改为把索引保存到文件中。这种把指针转换成索引、再把索引转换回指针的过程称为 marshmallowing。保存图书馆信息时，我们把整个保存过程分成若干步骤：

\begin{itemize}
\item \textbf{保存图书清单}：此时只保存作者与书名。
\item \textbf{保存顾客清单}：此时只保存姓名与地址。
\item \textbf{对每本图书}：保存其借阅者（如果已被借出）以及（可能为空的）预约列表。我们保存的是顾客列表的索引，而不是指向顾客的指针。
\item \textbf{对每位顾客}：保存其借阅与预约的图书集合。我们保存的是图书列表的索引，而不是指向图书的指针。
\end{itemize}

\subsection{将图书馆信息保存到文件}

\texttt{save} 方法打开文件，如果文件成功打开，就写出图书馆中的图书和顾客：

\begin{codebox}
void Library::save() {
  ofstream outStream(s_binaryPath);
\end{codebox}

\subsubsection{写入图书对象}

我们保存图书对象。这里只通过对每个 \texttt{Book} 对象调用 \texttt{write} 来保存书的作者与书名，暂不保存可能存在的借阅者和预约列表。

我们首先把列表中的图书数量写入文件：

\begin{codebox}
  if (outStream) {
    int bookPtrListSize = m_bookPtrList.size();
    outStream.write((char*) &bookPtrListSize,
                    sizeof bookPtrListSize);
\end{codebox}

然后，通过对每个 \texttt{Book} 对象指针调用 \texttt{write}，把每本书的信息写入文件：

\begin{codebox}
    for (const Book* bookPtr : m_bookPtrList) {
      bookPtr->write(outStream);
    }
  }
\end{codebox}

\subsubsection{写入顾客对象}

我们保存顾客对象。与前面的图书情况类似，这里只通过对每个 \texttt{Customer} 对象调用 \texttt{write} 来保存顾客的姓名与地址，暂不保存借阅与预约的图书集合。

与前面的图书情况相同，我们首先把列表中的顾客数量写入文件：

\begin{codebox}
  int customerPtrListSize = m_customerPtrList.size();
  outStream.write((char*) &customerPtrListSize,
                  sizeof customerPtrListSize);
\end{codebox}

然后，通过对每个 \texttt{Customer} 对象指针调用 \texttt{write} 方法，把每位顾客的信息写入文件：

\begin{codebox}
  for (const Customer* customerPtr : m_customerPtrList) {
    customerPtr->write(outStream);
  }
\end{codebox}

\subsubsection{写入借阅者索引}

对每个 \texttt{Book} 对象，如果这本书已被借出，我们就查找并保存相应 \texttt{Customer} 的索引，而不是保存指向该对象的指针：

\begin{codebox}
  for (const Book* bookPtr : m_bookPtrList) {
    const Customer* borrowerPtr = bookPtr->borrowerPtr();
\end{codebox}

对每本书，我们首先检查它是否已被借出。如果已被借出，就向文件写入值 \texttt{true}，表示这本书处于被借出状态：

\begin{codebox}
    if (borrowerPtr != nullptr) {
      bool borrowed = true;
      outStream.write((char*) &borrowed, sizeof borrowed);
\end{codebox}

然后，我们在顾客指针列表中查找借走这本书的顾客的索引，并把该索引写入文件：

\begin{codebox}
      int loanIndex = lookupCustomerIndex(borrowerPtr);
      outStream.write((char*) &loanIndex, sizeof loanIndex);
    }
\end{codebox}

如果这本书未被借出，我们只需向文件写入值 \texttt{false}，表示这本书没有被借阅：

\begin{codebox}
    else {
      bool borrowed = false;
      outStream.write((char*) &borrowed, sizeof borrowed);
    }
  }
\end{codebox}

\subsubsection{写入预约索引}

由于一本书可以同时被多位顾客预约，我们遍历预约列表，并保存列表中每位顾客的索引：

\begin{codebox}
    const list<Customer*>& reservationPtrList =
        bookPtr->reservationPtrList();
\end{codebox}

对每本书，我们首先把该书的预约数量写入文件：

\begin{codebox}
    int reserveSetSize = reservationPtrList.size();
    outStream.write((char*) &reserveSetSize,
                    sizeof reserveSetSize);
\end{codebox}

然后遍历预约列表，对每一项预约，查找并写出预约这本书的每位顾客的索引：

\begin{codebox}
    for (const Customer* customerPtr : reservationPtrList) {
      int customerIndex = lookupCustomerIndex(customerPtr);
      outStream.write((char*) &customerIndex,
                      sizeof customerIndex);
    }
  }
\end{codebox}

\subsubsection{写入借阅图书索引}

对每位顾客，我们保存其借阅的各本书的索引：先保存借阅列表的大小，再保存各本书的索引：

\begin{codebox}
  for (const Customer* customerPtr : m_customerPtrList) {
    const set<Book*>& loanPtrSet = customerPtr->loanPtrSet();
\end{codebox}

对每位顾客，我们首先把借阅数量写入文件：

\begin{codebox}
    int loanPtrSetSize = loanPtrSet.size();
    outStream.write((char*) &loanPtrSetSize,
                    sizeof loanPtrSetSize);
\end{codebox}

然后遍历借阅集合，对每一次借阅，查找并把每本书的索引写入文件：

\begin{codebox}
    for (const Book* bookPtr : loanPtrSet) {
      int bookIndex = lookupBookIndex(bookPtr);
      outStream.write((char*) &bookIndex,
                      sizeof bookIndex);
    }
  }
\end{codebox}

\subsubsection{写入预约图书索引}

同样，对每位顾客，我们保存其预约的各本书的索引：先保存预约列表的大小，再保存其所预约各本书的索引：

\begin{codebox}
    const set<Book*>& reservationPtrSet =
        customerPtr->reservationPtrSet();
\end{codebox}

对每位顾客，我们首先把预约的图书数量写入文件：

\begin{codebox}
    int reservationPtrSetSize = reservationPtrSet.size();
    outStream.write((char*) &reservationPtrSetSize,
                    sizeof reservationPtrSetSize);
\end{codebox}

然后遍历预约集合，对每一项预约，查找并把每本书的索引写入文件：

\begin{codebox}
    for (const Book* bookPtr : reservationPtrSet) {
      int bookIndex = lookupBookIndex(bookPtr);
      outStream.write((char*) &bookIndex,
                      sizeof bookIndex);
    }
  }
}
\end{codebox}

\subsection{从文件加载图书馆信息}

加载文件时，我们按照与保存文件时相同的方式来进行：

\begin{codebox}
void Library::load() {
  ifstream inStream(s_binaryPath);
\end{codebox}

\subsubsection{读取图书对象}

我们先读取图书列表的大小，再读取各本书本身。记住，到目前为止我们只读取了书的作者与书名：

\begin{codebox}
  if (inStream) {
    int bookPtrListSize;
\end{codebox}

我们首先读取图书数量：

\begin{codebox}
    inStream.read((char*) &bookPtrListSize,
                  sizeof bookPtrListSize);
\end{codebox}

然后读取各本书本身。对每一本书，我们动态分配一个 \texttt{Book} 对象，通过对该指针调用 \texttt{read} 来读取其信息，并把该指针加入图书指针列表：

\begin{codebox}
    for (int count = 0; count < bookPtrListSize; ++count) {
      Book *bookPtr = new Book();
      assert(bookPtr != nullptr);
      bookPtr->read(inStream);
      m_bookPtrList.push_back(bookPtr);
    }
  }
\end{codebox}

\subsubsection{读取顾客对象}

同样，我们先读取顾客列表的大小，再读取各位顾客本身。到目前为止，我们只读取了顾客的姓名与地址：

\begin{codebox}
  int customerPtrListSize;
\end{codebox}

我们首先读取顾客数量：

\begin{codebox}
  inStream.read((char*) &customerPtrListSize,
                sizeof customerPtrListSize);
\end{codebox}

然后读取各位顾客本身。对每一位顾客，我们动态分配一个 \texttt{Customer} 对象，通过对该指针调用 \texttt{read} 来读取其信息，并把该指针加入顾客指针列表：

\begin{codebox}
  for (int count = 0; count < customerPtrListSize; ++count) {
    Customer *customerPtr = new Customer();
    assert(customerPtr != nullptr);
    customerPtr->read(inStream);
    m_customerPtrList.push_back(customerPtr);
  }
\end{codebox}

\subsubsection{读取借阅者索引}

对每本书，我们读取借阅过它的顾客（如果有的话），以及预约这本书的顾客列表：

\begin{codebox}
  for (Book* bookPtr : m_bookPtrList) {
    bool borrowed;
    inStream.read((char*) &borrowed, sizeof borrowed);
\end{codebox}

如果 borrowed 为 \texttt{true}，说明这本书已被借出。此时我们读取该顾客的索引，然后查找相应 \texttt{Customer} 对象的指针，并把它登记为这本书的借阅者：

\begin{codebox}
    if (borrowed) {
      int loanIndex;
      inStream.read((char*) &loanIndex, sizeof loanIndex);
      bookPtr->borrowerPtr() = lookupCustomerPtr(loanIndex);
    }
\end{codebox}

如果 borrowed 为 \texttt{false}，说明这本书没有被借出。此时我们把指向借阅者顾客的指针设为 \texttt{nullptr}：

\begin{codebox}
    else {
      bookPtr->borrowerPtr() = nullptr;
    }
  }
\end{codebox}

\subsubsection{读取预约索引}

对每本书，我们还要读取预约列表：先读取列表的大小，再读取各条顾客索引：

\begin{codebox}
    list<Customer*>& reservationPtrList =
        bookPtr->reservationPtrList();
    int reservationPtrListSize;
\end{codebox}

我们首先读取这本书的预约数量：

\begin{codebox}
    inStream.read((char*) &reservationPtrListSize,
                  sizeof reservationPtrListSize);
\end{codebox}

对每一项预约，我们读取顾客的索引，并调用 \texttt{lookupCustomerPtr} 获得相应 \texttt{Customer} 对象的指针，然后把它加入这本书的预约指针列表：

\begin{codebox}
    for (int count = 0; count < reservationPtrListSize;
         ++count) {
      int customerIndex;
      inStream.read((char*) &customerIndex,
                    sizeof customerIndex);
      Customer* customerPtr =
          lookupCustomerPtr(customerIndex);
      reservationPtrList.push_back(customerPtr);
    }
  }
\end{codebox}

\subsubsection{读取借阅图书索引}

对每位顾客，我们读取其借阅的图书集合：

\begin{codebox}
  for (Customer* customerPtr : m_customerPtrList) {
    set<Book*>& loanPtrSet = customerPtr->loanPtrSet();
    int loanPtrSetSize = loanPtrSet.size();
\end{codebox}

我们首先读取借阅列表的大小：

\begin{codebox}
    inStream.read((char*) &loanPtrSetSize,
                  sizeof loanPtrSetSize);
\end{codebox}

对每一次借阅，我们读取该书的索引，并调用 \texttt{lookupBookPtr} 获得相应 \texttt{Book} 对象的指针，然后把它加入借阅指针集合：

\begin{codebox}
    for (int count = 0; count < loanPtrSetSize; ++count) {
      int bookIndex;
      inStream.read((char*) &bookIndex, sizeof bookIndex);
      Book* bookPtr = lookupBookPtr(bookIndex);
      loanPtrSet.insert(bookPtr);
    }
  }
\end{codebox}

\subsubsection{读取预约图书索引}

同样，对每位顾客，我们读取其预约的图书集合：

\begin{codebox}
    set<Book*>& reservationPtrSet =
        customerPtr->reservationPtrSet();
\end{codebox}

我们首先读取预约列表的大小：

\begin{codebox}
    int reservationPtrSetSize = reservationPtrSet.size();
    inStream.read((char*) &reservationPtrSetSize,
                  sizeof reservationPtrSetSize);
\end{codebox}

对每一项预约，我们读取该书的索引，并调用 \texttt{lookupBookPtr} 获得相应 \texttt{Book} 对象的指针，然后把它加入预约指针集合：

\begin{codebox}
    for (int count = 0; count < reservationPtrSetSize;
         ++count) {
      int bookIndex;
      inStream.read((char*) &bookIndex, sizeof bookIndex);
      Book* bookPtr = lookupBookPtr(bookIndex);
      reservationPtrSet.insert(bookPtr);
    }
  }
}
\end{codebox}

\subsection{释放内存}

由于我们向列表中添加的是动态分配的 \texttt{Book} 和 \texttt{Customer} 对象，因此需要在程序执行结束时释放它们。析构函数遍历图书指针列表和顾客指针列表，释放所有图书指针和顾客指针：

\begin{codebox}
Library::~Library() {
  for (const Book* bookPtr : m_bookPtrList) {
    delete bookPtr;
  }
  for (const Customer* customerPtr : m_customerPtrList) {
    delete customerPtr;
  }
}
\end{codebox}

\subsection{main 函数}

与第 3 章「构建图书馆管理系统」类似，main 函数只是创建一个 \texttt{Library} 对象：

\medskip\noindent{\ttfamily\bfseries Main.cpp}

\begin{codebox}
#include <Set>
#include <Map>
#include <String>
#include <FStream>
#include <IOStream>
using namespace std;

#include "Book.h"
#include "Customer.h"
#include "Library.h"

void main() {
  Library();
}
\end{codebox}

\section{小结}

本章我们构建了一个与第 3 章「构建图书馆管理系统」类似的图书馆管理系统。不过，我们省去了所有整数编号，改用指针。这样做的好处是可以更直接地登记借阅与预约，但代价是把它们保存到文件以及从文件加载变得更加困难。

在第 5 章「Qt 图形应用程序」中，我们将研究图形应用程序。
